% ============================================================================
% Fichier  : partie4/P04_C11_fondements_custody.tex
% Titre    : Fondements du Forensique et Chaîne de Custody
% Auteur   : MINKA MI NGUIDJOÏ Thierry Emmanuel
% ============================================================================
% Sources : ISO/IEC 27037 ; RFC 3227 ; NIST SP 800-86. Opération FARO : fictif.
% ============================================================================

\chapter{Fondements du Forensique et Chaîne de Custody}
\label{chap:fondements}
\chaptermark{Fondements du forensique}

\epigraph{%
    « L'enquêteur ne cherche pas à avoir raison~;\\
    il cherche à pouvoir le prouver. »%
}{Principe du cours}

\niveauChapitre{Fondamental}{Partie~III.}

\begin{boxobjectifs}
\begin{itemize}
    \item Définir l'analyse forensique et ses principes.
    \item Comprendre l'ordre de volatilité et la copie vérifiée.
    \item Tenir une chaîne de custody.
    \item Qualifier un artefact.
    \item Préserver l'opposabilité à chaque étape.
\end{itemize}
\end{boxobjectifs}

\bigskip

% ============================================================================
\section*{Question directrice : par où commence une enquête propre sur FS-01~?}
% ============================================================================

Devant un serveur compromis, la tentation est d'« aller voir ». Le forensique
impose l'inverse~: figer, copier, documenter, puis analyser une copie. Ce
chapitre pose les fondements et la chaîne de custody.

\medskip
\textbf{Plan du chapitre.}
\begin{enumerate}
    \item Qu'est-ce que le forensique (section~\ref{sec:c11-definir}).
    \item L'ordre de volatilité (section~\ref{sec:c11-volatilite}).
    \item La copie vérifiée (section~\ref{sec:c11-copie}).
    \item La chaîne de custody (section~\ref{sec:c11-custody}).
    \item Qualifier un artefact (section~\ref{sec:c11-artefact}).
\end{enumerate}

\separateur

% ============================================================================
\section{Qu'est-ce que le Forensique}
\label{sec:c11-definir}
% ============================================================================

\begin{boxdef}
\termecle{Analyse forensique}~: recherche méthodique, documentée et opposable
des traces laissées par un événement sur les systèmes, afin d'établir ce qui
s'est passé, comment, et par qui, dans le respect de l'intégrité des preuves.
\end{boxdef}

% ============================================================================
\section{L'Ordre de Volatilité}
\label{sec:c11-volatilite}
% ============================================================================

\begin{boxprocedure}
\textbf{Acquérir selon l'ordre de volatilité (RFC 3227)}
\begin{enumerate}
    \item Mémoire vive et état du système (processus, connexions).
    \item Fichiers temporaires et données de session.
    \item Disque (systèmes de fichiers).
    \item Journaux distants et configuration.
    \item Supports d'archives.
\end{enumerate}
\end{boxprocedure}

On capture d'abord ce qui disparaît le plus vite. Un redémarrage détruit la
mémoire~: d'où la règle « ne rien toucher avant d'avoir figé »
(Chapitre~\ref{chap:preuve}).

% ============================================================================
\section{La Copie Vérifiée}
\label{sec:c11-copie}
% ============================================================================

On n'analyse jamais l'original~: on en fait une \termecle{copie bit à bit},
dont on calcule un \termecle{condensat} (empreinte). Le même condensat sur
l'original et la copie prouve que la copie est fidèle et que l'original n'a
pas changé. Toute l'analyse se fait sur la copie.

% ============================================================================
\section{La Chaîne de Custody}
\label{sec:c11-custody}
% ============================================================================

\begin{boxprocedure}
\textbf{Tenir la chaîne de custody}
\begin{enumerate}
    \item Identifier chaque pièce (numéro, description, condensat).
    \item Consigner chaque transfert~: qui, à qui, quand, pourquoi.
    \item Protéger physiquement et logiquement les pièces.
    \item Ne jamais laisser de trou dans la continuité de la détention.
    \item Joindre la chaîne au dossier de preuve.
\end{enumerate}
\end{boxprocedure}

\begin{boxCRO}
\textbf{La chaîne de custody sert l'opposabilité.} Une pièce dont la
détention présente un trou devient contestable~: \cro{0}{0}{-}. La continuité
documentée est ce qui rend la preuve opposable devant un tiers.
\end{boxCRO}

% ============================================================================
\section{Qualifier un Artefact}
\label{sec:c11-artefact}
% ============================================================================

\begin{boxobjet}{Journal de connexion de FS-01}
\textbf{Nature~:} journal d'authentification extrait de la copie du disque de
\texttt{FS-01}. \textbf{Source~:} copie vérifiée (condensat au dossier).
\textbf{Pertinence~:} contient les connexions du compte de service suspect.
\textbf{Fiabilité~:} intègre (condensat), horodatage à recaler (écart connu
de 40~min). \textbf{Opposabilité~:} chaîne de custody continue.
\end{boxobjet}

\begin{boxmetier}{Analyste forensique}
\textbf{Documenter est enquêter.} Pour l'analyste, la documentation n'est pas
une formalité postérieure~: c'est une partie de l'acte d'enquête. Un geste
non documenté est un geste perdu.
\end{boxmetier}

\separateur

\begin{boxresume}
\textbf{Ce qu'il faut retenir de ce chapitre~:}
\begin{itemize}
    \item \textbf{Le forensique} établit les faits de façon opposable.
    \item \textbf{L'ordre de volatilité}~: du plus fugace au plus stable.
    \item \textbf{La copie vérifiée} par condensat remplace l'original.
    \item \textbf{La chaîne de custody} assure la continuité de la détention.
    \item \textbf{Un artefact} se qualifie (nature, source, pertinence,
          fiabilité, opposabilité).
\end{itemize}
\end{boxresume}

\begin{boxinfo}
\textbf{Points de contrôle de l'élève.} Je sais~: définir le forensique~;
énoncer l'ordre de volatilité~; faire une copie vérifiée~; tenir une chaîne de
custody~; qualifier un artefact.
\end{boxinfo}

\bigskip

\begin{boxexo}[title={\ding{45}\quad Exercice~11.1. Volatilité \niveaubase}]
Classez du plus volatil au plus stable~: disque, mémoire vive, archives.
\end{boxexo}

\begin{boxexo}[title={\ding{45}\quad Exercice~11.2. Copie \niveaubase}]
Pourquoi analyse-t-on une copie et non l'original~?
\end{boxexo}

\begin{boxexo}[title={\ding{45}\quad Exercice~11.3. Custody \niveaubase}]
Qu'est-ce qu'un « trou » dans la chaîne de custody~?
\end{boxexo}

\begin{boxexo}[title={\ding{45}\quad Exercice~11.4. Condensat \niveaumoyen}]
Comment un condensat prouve-t-il qu'une copie est fidèle~?
\end{boxexo}

\begin{boxexo}[title={\ding{45}\quad Exercice~11.5. CRO \niveaumoyen}]
Pourquoi un trou de custody dégrade-t-il l'opposabilité~?
\end{boxexo}

\begin{boxexo}[title={\ding{45}\quad Exercice~11.6. Artefact \niveaumoyen}]
Qualifiez un artefact de votre choix de l'opération FARO (cinq rubriques).
\end{boxexo}

\begin{boxexo}[title={\ding{45}\quad Exercice~11.7. Horloge \niveaumoyen}]
Comment traiter l'écart d'horloge de 40~min dans l'analyse~?
\end{boxexo}

\begin{boxexo}[title={\ding{45}\quad Exercice~11.8. Documentation \niveaumoyen}]
Pourquoi un geste non documenté est-il « perdu »~?
\end{boxexo}

\begin{boxexo}[title={\ding{45}\quad Exercice~11.9. Procédure d'acquisition \niveauexpert}]
Rédigez la procédure d'acquisition du serveur \texttt{FS-01}, dans l'ordre de
volatilité.
\end{boxexo}

\begin{boxexo}[title={\ding{45}\quad Exercice~11.10. Formulaire de custody \niveauexpert}]
Concevez le formulaire de chaîne de custody de LiZbiZ.
\end{boxexo}

\bigskip

\begin{boxinfo}
\textbf{Pour aller plus loin.}
\begin{itemize}
    \item ISO/IEC 27037~\cite{iso27037}~; RFC~3227~\cite{rfc3227}~;
          NIST SP~800-86~\cite{nist80086}.
\end{itemize}
\end{boxinfo}

\begin{boxinfo}
\textbf{Transition.} Le premier terrain d'enquête est le disque~: le
Chapitre~\ref{chap:systeme} traite du forensique système.
\end{boxinfo}
